\documentclass[11pt]{article}

\usepackage[margin=1in]{geometry}
\usepackage{amsmath,amssymb,amsthm,mathtools}
\usepackage{bm}
\usepackage{enumitem}
\usepackage{microtype}
\usepackage{hyperref}
\newtheorem{theorem}{Theorem}
\newtheorem{lemma}{Lemma}
\newtheorem{remark}{Remark}

\newcommand{\R}{\mathbb{R}}
\newcommand{\op}{\mathrm{op}}
\newcommand{\FW}{\mathsf{FW}}
\newcommand{\Sep}{\mathsf{Sep}}
\newcommand{\polylog}{\operatorname{polylog}}
\newcommand{\eps}{\varepsilon}

\title{Frank--Wolfe Complexity Separation Between\\
Frobenius and Operator-Norm Geometries}
\author{}
\date{}

\begin{document}

\maketitle

\begin{abstract}
We compare the standard smoothness--radius factors governing the
Frank--Wolfe convergence bounds under the Frobenius and operator-norm
geometries. For the Gaussian associative-memory softmax loss with
Zipf probabilities
\[
p_i=\frac{i^{-\alpha}}{H_{N,\alpha}},
\qquad
H_{N,\alpha}:=\sum_{j=1}^N j^{-\alpha},
\]
we study
\[
\FW_F(\eps):=\beta_F R_F^2(\eps),
\qquad
\FW_{\op}(\eps):=\beta_{\op}R_{\op}^2(\eps).
\]
We assume
\[
N=d^\kappa,
\qquad
0<\kappa<2,
\qquad
\alpha\ge0
\]
with $\kappa$ and $\alpha$ fixed. In the accuracy range where the
available Frobenius-radius lower bound is nontrivial, we prove a
polynomial separation between the two Frank--Wolfe complexity factors.
The dependence on the target loss $\eps$ cancels up to logarithmic
factors, uniformly over the entire admissible accuracy range.
\end{abstract}

\section{Setting and interpretation}

Let
\[
L(W) = \sum_{i=1}^N p_i \log\left( 1+\sum_{j\ne i} \exp\bigl((u_j-u_i)^\top Wv_i\bigr) \right),
\]
where the input and output embeddings are independent normalized
Gaussian vectors in $\R^d$.

For $X\in\{F,\op\}$, define
\[
R_X(\eps) := \inf\bigl\{||W||_X:L(W)\le\eps\bigr\}.
\]
Let $\beta_X$ denote the global smoothness constant of $L$ with
respect to $||\cdot||_X$.

For a $\beta_X$-smooth objective over an $X$-norm ball of radius $R$,
the Frank--Wolfe curvature constant is bounded, up to a universal
constant, by
\[
C_L\lesssim \beta_XR^2.
\]
Consequently, the geometry-dependent factor in the standard
Frank--Wolfe convergence estimate is
\[
\FW_X(\eps):=\beta_XR_X^2(\eps).
\]
This note compares these standard smoothness--radius factors. It does
not claim a minimax lower bound on the actual number of iterations of
every possible algorithm.

Throughout, $\widetilde O$, $\widetilde\Omega$, and
$\widetilde\Theta$ suppress powers of
\[
\log d,\qquad \log N,\qquad \log(1/\delta),
\]
but do not suppress polynomial powers of $d$ or $N$. Constants may
depend on the fixed Zipf exponent $\alpha$.

We assume
\[
\delta^{-1}=d^{O(1)}.
\]
The operator-radius construction requires
\[
N\le c_0\frac{d^2}{\log d}.
\]
Since $N=d^\kappa$ with fixed $\kappa<2$,
\[
\frac{N\log d}{d^2} = \frac{\log d}{d^{2-\kappa}} \longrightarrow0,
\]
so this condition holds for all sufficiently large $d$.

\begin{theorem}[Main Complexity Separation]
\label{thm:main}
Consider the Gaussian associative-memory softmax loss with Zipf probabilities $p_i \propto i^{-\alpha}$ for a fixed exponent $\alpha \ge 0$. Assume the number of facts is $N=d^\kappa$ for a fixed $0<\kappa<2$, and let the failure probability be $\delta^{-1}=d^{O(1)}$.

There exist constants $c_\alpha > 0$ such that for any target accuracy $\eps$ in the range
\[
0 < \eps \le c_\alpha
\begin{cases}
1, & 0 \le \alpha < 1, \\[1mm]
(\log N)^{-1}, & \alpha = 1, \\[1mm]
N^{1-\alpha}, & \alpha > 1,
\end{cases}
\]
the ratio of the Frank--Wolfe smoothness--radius factors satisfies, with probability at least $1-\delta$,
\[
\boxed{
\frac{\FW_F(\eps)}{\FW_{\op}(\eps)} 
:= \frac{\beta_F R_F^2(\eps)} {\beta_{\op}R_{\op}^2(\eps)} =
\begin{cases}
\widetilde\Omega_\alpha(d^{\kappa\alpha}), & 0<\kappa\le1,\quad 0\le\alpha<1, \\[2mm]
\widetilde\Omega_\alpha(d^\kappa), & 0<\kappa\le1,\quad \alpha\ge1, \\[2mm]
\widetilde\Omega_\alpha\left( d^{(\kappa-1)(1-\alpha)}+d^\alpha \right), & 1<\kappa<2,\quad 0\le\alpha<1, \\[2mm]
\widetilde\Omega_\alpha(d), & 1<\kappa<2,\quad \alpha\ge1.
\end{cases}
}
\]
\end{theorem}

\section{Input estimates}

Set
\[
m:=\min\{d,N\},
\qquad
P_m:=\sum_{i=1}^m p_i =\frac{H_{m,\alpha}}{H_{N,\alpha}}.
\]

The smoothness estimates give, with high probability,
\begin{equation}
\beta_F \asymp \frac1d+p_1 = \frac1d+\frac1{H_{N,\alpha}},
\label{eq:beta-F-general}
\end{equation}
and
\begin{equation}
\beta_{\op} \lesssim_\alpha (\log d)^2 \left( \frac1d+P_m \right).
\label{eq:beta-op-general}
\end{equation}
The power $(\log d)^2$ in \eqref{eq:beta-op-general} is a convenient
uniform upper bound. For fixed $\alpha>0$, the available estimate can
be sharpened by one logarithmic factor, but this does not affect any
polynomial separation derived below.

The operator-radius construction gives, for every $0<\eps\le1$,
\begin{equation}
R_{\op}^2(\eps) \lesssim \left(1+\frac Nd\right) \log^2\left(\frac{2N}{\eps}\right).
\label{eq:Rop}
\end{equation}

For the Frobenius radius, introduce the natural probability scale
\begin{equation}
q_\alpha(N) :=
\begin{cases}
1, & 0\le\alpha<1,\\[1mm]
(\log N)^{-1}, & \alpha=1,\\[1mm]
N^{1-\alpha}, & \alpha>1.
\end{cases}
\label{eq:q-alpha}
\end{equation}
After changing constants depending only on $\alpha$, the Frobenius
lower bounds can be written uniformly as
\begin{equation}
R_F^2(\eps) \gtrsim_\alpha \frac{N}{\log d} \log^2\left( \frac{c_\alpha q_\alpha(N)}{\eps} \right),
\qquad
0<\eps\le c_\alpha q_\alpha(N).
\label{eq:RF-unified}
\end{equation}

More explicitly,
\begin{equation}
R_F^2(\eps) \gtrsim_\alpha
\begin{cases}
\displaystyle
\frac{N}{\log d} \log^2\left(\frac{c_\alpha}{\eps}\right),
& 0\le\alpha<1, \quad 0<\eps\le c_\alpha, \\[4mm]
\displaystyle
\frac{N}{\log d} \log^2\left(\frac{c}{\eps\log N}\right),
& \alpha=1, \quad 0<\eps\le \frac{c}{\log N}, \\[4mm]
\displaystyle
\frac{N}{\log d} \log^2\left( \frac{c_\alpha}{\eps N^{\alpha-1}} \right),
& \alpha>1, \quad 0<\eps\le c_\alpha N^{1-\alpha}.
\end{cases}
\label{eq:RF-cases}
\end{equation}

\section{Why the accuracy dependence disappears in the ratio}

The logarithms in \eqref{eq:Rop} and \eqref{eq:RF-unified} do not
literally coincide. Nevertheless, their ratio is bounded below by an
inverse polylogarithmic factor uniformly over the entire relevant
accuracy range.

\begin{lemma}[Uniform comparison of the accuracy logarithms]
\label{lem:epsilon-cancel}
After decreasing $c_\alpha$ if necessary, suppose
\[
0<\eps\le c_\alpha q_\alpha(N)
\]
and define
\[
\ell_F(\eps) := \log\left( \frac{c_\alpha q_\alpha(N)}{\eps} \right),
\qquad
\ell_{\op}(\eps) := \log\left(\frac{2N}{\eps}\right).
\]
Then
\[
\frac{\ell_F^2(\eps)}{\ell_{\op}^2(\eps)} \gtrsim_\alpha \frac1{(\log d)^2}.
\]
In particular, no polynomial dependence on $\eps$ remains in the
complexity ratio.
\end{lemma}

\begin{proof}
By shrinking $c_\alpha$, we may assume
\[
\ell_F(\eps)\ge1.
\]
Moreover,
\begin{align}
\ell_{\op}(\eps)
&= \log\left(\frac{2N}{\eps}\right) \nonumber \\
&= \log\left( \frac{c_\alpha q_\alpha(N)}{\eps} \right) + \log\left( \frac{2N}{c_\alpha q_\alpha(N)} \right) \nonumber \\
&= \ell_F(\eps) + \log\left( \frac{2N}{c_\alpha q_\alpha(N)} \right).
\label{eq:log-decomposition}
\end{align}
For fixed $\alpha$,
\[
\frac{N}{q_\alpha(N)} =
\begin{cases}
N, & 0\le\alpha<1,\\
N\log N, & \alpha=1,\\
N^\alpha, & \alpha>1.
\end{cases}
\]
Since $N=d^\kappa$ and $\kappa$ is fixed,
\[
\log\left( \frac{2N}{c_\alpha q_\alpha(N)} \right) \lesssim_{\alpha,\kappa} \log d.
\]
Using $\ell_F(\eps)\ge1$ in \eqref{eq:log-decomposition},
\[
\ell_{\op}(\eps) \le \ell_F(\eps)+C_{\alpha,\kappa}\log d \le \bigl(1+C_{\alpha,\kappa}\log d\bigr) \ell_F(\eps).
\]
Therefore,
\[
\frac{\ell_F^2(\eps)}{\ell_{\op}^2(\eps)} \ge \frac1{ \bigl(1+C_{\alpha,\kappa}\log d\bigr)^2 } \gtrsim_{\alpha,\kappa} \frac1{(\log d)^2}.
\]
\end{proof}

\begin{remark}
The restriction
\[
\eps\lesssim_\alpha q_\alpha(N)
\]
is essential for the present Frobenius-radius lower bound. In
particular, when $\alpha>1$, the relevant accuracy range is
\[
\eps\lesssim_\alpha N^{1-\alpha}.
\]
At larger accuracies, the population loss may ignore a block
containing $\Theta(N)$ low-probability facts, so the current
$K\asymp N$ Frobenius lower bound need not remain nontrivial.
\end{remark}

\section{A unified Frank--Wolfe separation formula}

Define
\[
\Sep(\eps) := \frac{\FW_F(\eps)}{\FW_{\op}(\eps)} = \frac{\beta_FR_F^2(\eps)} {\beta_{\op}R_{\op}^2(\eps)}.
\]

\begin{theorem}[Unified complexity separation]
\label{thm:unified}
Let $\alpha\ge0$ and $N=d^\kappa$ with fixed $0<\kappa<2$.
For
\[
0<\eps\le c_\alpha q_\alpha(N),
\]
with high probability,
\begin{equation}
\Sep(\eps) \gtrsim_\alpha \frac1{(\log d)^5} \min\{N,d\} \frac{d^{-1}+p_1} {d^{-1}+P_{\min\{N,d\}}}.
\label{eq:explicit-polylog-separation}
\end{equation}
Equivalently,
\begin{equation}
\boxed{
\Sep(\eps) = \widetilde\Omega_\alpha\left( \min\{N,d\} \frac{d^{-1}+p_1} {d^{-1}+P_{\min\{N,d\}}} \right).
}
\label{eq:master-separation}
\end{equation}
\end{theorem}

\begin{proof}
Combining \eqref{eq:beta-F-general} and
\eqref{eq:RF-unified},
\[
\FW_F(\eps) \gtrsim_\alpha \left(\frac1d+p_1\right) \frac{N}{\log d} \ell_F^2(\eps).
\]
Combining \eqref{eq:beta-op-general} and \eqref{eq:Rop},
\[
\FW_{\op}(\eps) \lesssim_\alpha (\log d)^2 \left(\frac1d+P_m\right) \left(1+\frac Nd\right) \ell_{\op}^2(\eps).
\]
Thus
\begin{align}
\Sep(\eps) &\gtrsim_\alpha \frac{N}{1+N/d} \frac{d^{-1}+p_1}{d^{-1}+P_m} \frac1{(\log d)^3} \frac{\ell_F^2(\eps)}{\ell_{\op}^2(\eps)}.
\label{eq:ratio-before-log}
\end{align}
Lemma~\ref{lem:epsilon-cancel} gives
\[
\frac{\ell_F^2(\eps)}{\ell_{\op}^2(\eps)} \gtrsim_\alpha \frac1{(\log d)^2}.
\]
Finally,
\[
\frac{N}{1+N/d} = \frac{Nd}{N+d} \asymp \min\{N,d\}.
\]
Substitution into \eqref{eq:ratio-before-log} proves
\eqref{eq:explicit-polylog-separation} and
\eqref{eq:master-separation}.
\end{proof}

\section{Evaluation of the Zipf quantities}

The standard generalized-harmonic-number estimates are
\[
H_{n,\alpha} \asymp_\alpha
\begin{cases}
n^{1-\alpha}, & 0\le\alpha<1,\\
\log n, & \alpha=1,\\
1, & \alpha>1.
\end{cases}
\]
Consequently,
\begin{equation}
p_1 = \frac1{H_{N,\alpha}} \asymp_\alpha
\begin{cases}
N^{\alpha-1}, & 0\le\alpha<1,\\
(\log N)^{-1}, & \alpha=1,\\
1, & \alpha>1,
\end{cases}
\label{eq:p1}
\end{equation}
and
\begin{equation}
P_m = \frac{H_{m,\alpha}}{H_{N,\alpha}} \asymp_\alpha
\begin{cases}
\left(\dfrac mN\right)^{1-\alpha}, &0\le\alpha<1,\\[3mm]
\dfrac{\log m}{\log N}, &\alpha=1,\\[3mm]
1, &\alpha>1.
\end{cases}
\label{eq:Pm}
\end{equation}

We now evaluate \eqref{eq:master-separation} in all regimes.

\section{Regime I: \texorpdfstring{$N\le d$}{N less than d}}

Suppose
\[
0<\kappa\le1, \qquad N=d^\kappa\le d.
\]
Then
\[
m=N, \qquad P_m=P_N=1.
\]
Therefore,
\[
\Sep(\eps) = \widetilde\Omega_\alpha\left( N\left(\frac1d+p_1\right) \right).
\]

\subsection{\texorpdfstring{$0\le\alpha<1$}{alpha less than one}}

In this regime,
\[
p_1\asymp_\alpha N^{\alpha-1}.
\]
Since $N\le d$ and $0<1-\alpha\le1$,
\[
N^{\alpha-1} = \frac1{N^{1-\alpha}} \ge \frac1d.
\]
Hence
\[
\frac1d+p_1 \asymp_\alpha N^{\alpha-1},
\]
and thus
\begin{equation}
\boxed{
\Sep(\eps) = \widetilde\Omega_\alpha(d^{\kappa\alpha}),
\qquad
0\le\alpha<1,\quad 0<\kappa\le1.
}
\label{eq:sep-smallN-alpha-less1}
\end{equation}

At $\alpha=0$, this gives only
\[
\Sep(\eps)=\widetilde\Omega(1),
\]
so the present estimates do not yield a polynomial separation in the
uniform-frequency regime when $N\le d$.

For every fixed $\alpha>0$, however,
\[
d^{\kappa\alpha}\to\infty,
\]
and hence there is a polynomial advantage for the operator geometry.

\subsection{\texorpdfstring{$\alpha=1$}{alpha equals one}}

Here
\[
p_1\asymp\frac1{\log N},
\]
which dominates $1/d$. Therefore,
\[
\Sep(\eps) = \widetilde\Omega\left( \frac{N}{\log N} \right) = \widetilde\Omega(N).
\]
Since $N=d^\kappa$,
\begin{equation}
\boxed{
\Sep(\eps) = \widetilde\Omega(d^\kappa),
\qquad
\alpha=1,\quad 0<\kappa\le1.
}
\label{eq:sep-smallN-alpha1}
\end{equation}

\subsection{\texorpdfstring{$\alpha>1$}{alpha greater than one}}

In this case,
\[
p_1\asymp_\alpha1,
\]
so
\begin{equation}
\boxed{
\Sep(\eps) = \widetilde\Omega_\alpha(d^\kappa),
\qquad
\alpha>1,\quad 0<\kappa\le1.
}
\label{eq:sep-smallN-alpha-greater1}
\end{equation}

Combining all three cases,
\begin{equation}
\boxed{
\Sep(\eps) =
\begin{cases}
\widetilde\Omega_\alpha(d^{\kappa\alpha}), &0\le\alpha<1,\\[1mm]
\widetilde\Omega_\alpha(d^\kappa), &\alpha\ge1,
\end{cases}
\qquad
0<\kappa\le1.
}
\label{eq:sep-kappa-small-final}
\end{equation}
Equivalently,
\[
\Sep(\eps) = \widetilde\Omega_\alpha\left( d^{\kappa\min\{\alpha,1\}} \right), \qquad 0<\kappa\le1.
\]

\section{Regime II: \texorpdfstring{$d<N<d^2$}{d less than N less than d squared}}

Suppose
\[
1<\kappa<2, \qquad N=d^\kappa.
\]
Then
\[
m=d.
\]

\subsection{\texorpdfstring{$0\le\alpha<1$}{alpha less than one}}

We have
\[
p_1\asymp_\alpha N^{\alpha-1}, \qquad P_d \asymp_\alpha \left(\frac dN\right)^{1-\alpha}.
\]
Moreover, because $d<N<d^2$,
\[
\left(\frac dN\right)^{1-\alpha} \ge \frac dN > \frac1d.
\]
Therefore,
\[
\frac1d+P_d \asymp_\alpha \left(\frac dN\right)^{1-\alpha}.
\]
The master formula gives
\begin{align}
\Sep(\eps)
&= \widetilde\Omega_\alpha\left[ d \frac{d^{-1}+N^{\alpha-1}} {(d/N)^{1-\alpha}} \right] \nonumber \\
&= \widetilde\Omega_\alpha\left[ \frac{1}{(d/N)^{1-\alpha}} + \frac{dN^{\alpha-1}} {(d/N)^{1-\alpha}} \right] \nonumber \\
&= \widetilde\Omega_\alpha\left[ \left(\frac Nd\right)^{1-\alpha} +d^\alpha \right].
\label{eq:largeN-alpha-less1-Nd}
\end{align}
Thus
\begin{equation}
\boxed{
\Sep(\eps) = \widetilde\Omega_\alpha\left[ \left(\frac Nd\right)^{1-\alpha} +d^\alpha \right], \qquad d<N<d^2,\quad 0\le\alpha<1.
}
\label{eq:largeN-alpha-less1}
\end{equation}

Using $N=d^\kappa$,
\begin{equation}
\boxed{
\Sep(\eps) = \widetilde\Omega_\alpha\left[ d^{(\kappa-1)(1-\alpha)} +d^\alpha \right] = \widetilde\Omega_\alpha\left( d^{\max\{(\kappa-1)(1-\alpha), \alpha\}} \right).
}
\label{eq:largeN-kappa-alpha-less1}
\end{equation}

The two powers are equal when
\[
(\kappa-1)(1-\alpha)=\alpha.
\]
Solving gives the phase boundary
\begin{equation}
\alpha_\star = 1-\frac1\kappa.
\label{eq:alpha-star}
\end{equation}
Therefore,
\begin{equation}
\boxed{
\Sep(\eps) =
\begin{cases}
\widetilde\Omega_\alpha\left( d^{(\kappa-1)(1-\alpha)} \right), & 0\le\alpha\le 1-\dfrac1\kappa, \\[3mm]
\widetilde\Omega_\alpha(d^\alpha), & 1-\dfrac1\kappa\le\alpha<1.
\end{cases}
}
\label{eq:largeN-alpha-less1-piecewise}
\end{equation}

The two expressions agree at
$\alpha=1-1/\kappa$.

The origin of the two terms in
\eqref{eq:largeN-alpha-less1-Nd} is worth emphasizing:
\[
\left(\frac Nd\right)^{1-\alpha}
\]
comes from the diffuse $d^{-1}$ component of $\beta_F$, whereas
\[
d^\alpha
\]
comes from the largest-atom component
$p_1\asymp N^{\alpha-1}$.

In particular, at $\alpha=0$,
\[
\boxed{
\Sep(\eps) = \widetilde\Omega\left(d^{\kappa-1}\right).
}
\]
Thus, even for uniform frequencies, a polynomial separation emerges
as soon as $N>d$.

\subsection{\texorpdfstring{$\alpha=1$}{alpha equals one}}

For $N=d^\kappa$ with $\kappa>1$,
\[
p_1\asymp\frac1{\log N}, \qquad P_d \asymp \frac{\log d}{\log N} = \frac1\kappa \asymp1.
\]
Hence
\[
\Sep(\eps) = \widetilde\Omega\left( d\frac{1/\log N}{1} \right) = \widetilde\Omega\left( \frac d{\log N} \right).
\]
Since $\log N=\kappa\log d$,
\begin{equation}
\boxed{
\Sep(\eps) = \widetilde\Omega(d), \qquad \alpha=1,\quad 1<\kappa<2.
}
\label{eq:largeN-alpha1}
\end{equation}

\subsection{\texorpdfstring{$\alpha>1$}{alpha greater than one}}

For fixed $\alpha>1$,
\[
p_1\asymp_\alpha1, \qquad P_d\asymp_\alpha1.
\]
Therefore,
\begin{equation}
\boxed{
\Sep(\eps) = \widetilde\Omega_\alpha(d), \qquad \alpha>1,\quad 1<\kappa<2.
}
\label{eq:largeN-alpha-greater1}
\end{equation}

Combining the three cases,
\begin{equation}
\boxed{
\Sep(\eps) =
\begin{cases}
\widetilde\Omega_\alpha\left[ d^{(\kappa-1)(1-\alpha)}+d^\alpha \right], & 0\le\alpha<1, \\[2mm]
\widetilde\Omega_\alpha(d), & \alpha\ge1,
\end{cases}
\qquad
1<\kappa<2.
}
\label{eq:sep-kappa-large-final}
\end{equation}

Equivalently, the polynomial exponent can be written compactly as
\[
\Sep(\eps) = \widetilde\Omega_\alpha\left(d^{s(\kappa,\alpha)}\right),
\]
where
\begin{equation}
s(\kappa,\alpha) =
\begin{cases}
\kappa\min\{\alpha,1\}, &0<\kappa\le1, \\[2mm]
\min\left\{ 1, \max\left\{ \alpha, (\kappa-1)(1-\alpha) \right\} \right\}, &1<\kappa<2.
\end{cases}
\label{eq:separation-exponent}
\end{equation}

\section{Equivalent comparison of the two complexity factors}

The preceding results can also be understood by comparing the
polynomial parts of the numerator and denominator separately.

\subsection*{Case \(N\le d\)}

Since
\[
\beta_{\op}=\Theta(1), \qquad R_{\op}^2(\eps) \lesssim \log^2\left(\frac{2N}{\eps}\right),
\]
the operator complexity factor has no polynomial dependence on $d$ or
$N$:
\[
\FW_{\op}(\eps) \le \polylog(d,N,1/\eps).
\]
On the other hand,
\[
\FW_F(\eps) \gtrsim_\alpha \frac{N}{\log d} \left(\frac1d+p_1\right) \ell_F^2(\eps).
\]
Its polynomial part is therefore
\[
\begin{cases}
N^\alpha, & 0\le\alpha<1,\\
N, & \alpha\ge1,
\end{cases}
\]
up to logarithmic factors.

\subsection*{Case \(d<N<d^2\)}

For $0\le\alpha<1$,
\[
\beta_{\op} = \widetilde O_\alpha\left( \left(\frac dN\right)^{1-\alpha} \right), \qquad R_{\op}^2(\eps) \lesssim \frac Nd \log^2\left(\frac{2N}{\eps}\right).
\]
Thus the polynomial part of the operator complexity factor is
\[
\FW_{\op}(\eps) = \widetilde O_\alpha\left( \left(\frac Nd\right)^\alpha \right),
\]
whereas
\[
\FW_F(\eps) = \widetilde\Omega_\alpha\left( \frac Nd+N^\alpha \right)
\]
after matching the accuracy logarithms. Consequently,
\begin{align*}
\frac{\FW_F(\eps)}{\FW_{\op}(\eps)}
&= \widetilde\Omega_\alpha\left( \frac{N/d+N^\alpha}{(N/d)^\alpha} \right)\\
&= \widetilde\Omega_\alpha\left[ \left(\frac Nd\right)^{1-\alpha} +d^\alpha \right].
\end{align*}

For $\alpha\ge1$,
\[
\beta_{\op}=\Theta_\alpha(1), \qquad R_{\op}^2(\eps) \lesssim \frac Nd \log^2\left(\frac{2N}{\eps}\right),
\]
while the Frobenius polynomial factor is of order $N$.
The resulting separation is therefore of order $d$, up to
polylogarithmic factors.

\section{Final conclusion}

For all fixed $\alpha\ge0$ and fixed $0<\kappa<2$, in the accuracy
range
\[
0<\eps\lesssim_\alpha
\begin{cases}
1, & 0\le\alpha<1,\\
(\log N)^{-1}, & \alpha=1,\\
N^{1-\alpha}, & \alpha>1,
\end{cases}
\]
the Frank--Wolfe smoothness--radius factors satisfy
\begin{equation}
\boxed{
\frac{\beta_FR_F^2(\eps)} {\beta_{\op}R_{\op}^2(\eps)} =
\begin{cases}
\widetilde\Omega_\alpha(d^{\kappa\alpha}), & 0<\kappa\le1,\quad 0\le\alpha<1, \\[2mm]
\widetilde\Omega_\alpha(d^\kappa), & 0<\kappa\le1,\quad \alpha\ge1, \\[2mm]
\widetilde\Omega_\alpha\left( d^{(\kappa-1)(1-\alpha)}+d^\alpha \right), & 1<\kappa<2,\quad 0\le\alpha<1, \\[2mm]
\widetilde\Omega_\alpha(d), & 1<\kappa<2,\quad \alpha\ge1.
\end{cases}
}
\end{equation}

For $1<\kappa<2$ and $0\le\alpha<1$, the more explicit phase diagram is
\[
\boxed{
\frac{\beta_FR_F^2(\eps)} {\beta_{\op}R_{\op}^2(\eps)} =
\begin{cases}
\widetilde\Omega_\alpha\left( d^{(\kappa-1)(1-\alpha)} \right), & 0\le\alpha\le1-\dfrac1\kappa, \\[3mm]
\widetilde\Omega_\alpha(d^\alpha), & 1-\dfrac1\kappa\le\alpha<1.
\end{cases}
}
\]

Thus, except for the corner
\[
\alpha=0,\qquad N\le d,
\]
the available bounds yield a polynomial separation in favor of the
operator-norm geometry. The maximal separation is linear in the
ambient dimension:
\[
\Sep(\eps)=\widetilde\Omega_\alpha(d)
\]
throughout the heavy-head regime $\alpha\ge1$ when $N>d$.

\end{document}